\chapter{Training and the Run Contract}

\ChapterMeta{
This chapter explains the training stack and the run interface contract: stable artifact layout, strict config resolution, and resumable runs.
}{
Following the run interface contract keeps runs reproducible, makes parity checks easier, and keeps outputs predictable across machines and branches.
}{
It assumes run-interface-contract mode, such as \cmd{--run-id} or \cmd{--run-contract}, strict config keys, and the standard training dependencies.
}
\ChapterAudience{Read this chapter when you are running training jobs and want reproducible artifact layout, resume behavior, and tuning surfaces.}

\WorkflowOpening
{Run training so checkpoints, summaries, exports, resume state, and evaluation handoff files land in predictable locations.}
{Use this when YOLOZU owns the reference training run or when an external lane must emit a stable training summary and next-step handoff.}
{A normalized dataset root or wrapper, a YAML/JSON training config, and optional resume/export/parity settings.}
{\cmd{yolozu train ...} for the public route, or \cmd{python3 rtdetr_pose/tools/train_minimal.py ...} when using the low-level reference trainer.}
{\path{runs/<run_id>/checkpoints/}, \path{runs/<run_id>/reports/}, optional \path{exports/}, and \path{training_summary.json}.}
{Check that \path{config_resolved.yaml}, metrics JSONL, checkpoint paths, and optional export/parity reports exist and match the selected run id.}
{Native dataset layouts passed directly to the reference trainer, missing config keys, stale resume paths, backend-specific export failures, or unexpected training dependencies.}

\section{Model-family policy and reference training stack}
YOLOZU interface contracts for predictions, evaluation reports, and scenario outputs are model-family agnostic. The in-repo training reference implementation is the RT-DETR-style adapter under \path{rtdetr_pose/}. The critical part is the \textbf{run interface contract}: stable artifact paths and resumable training runs, independent of the model family that produced them.

\subsection{Dataset-format boundary for the reference trainer}
The RT-DETR reference trainer does not claim that it can ingest every native dataset layout directly at training time. Its expected input is a dataset root that \cmd{build_manifest(...)} can resolve: a YOLO-style root, a YOLO/Ultralytics \path{data.yaml}, or a \path{dataset.json} wrapper with \path{images_dir} and \path{labels_dir} pointing to image files and YOLO label files. The multi-format user convenience layer therefore lives one step earlier: \cmd{yolozu doctor import --dataset-from auto}, \cmd{yolozu import dataset --from auto}, and \cmd{yolozu migrate dataset --from auto} can auto-detect supported native sources such as COCO, COCO keypoints, Ultralytics/YOLO, and the supported semantic-segmentation roots, then materialize a normalized wrapper for training.

In other words, users do not need to manually classify the incoming dataset family, but the reference trainer itself still consumes a train-ready YOLOZU/YOLO-style result of that preflight/import step. Semantic-segmentation descriptors and SynthGen intake shards are import/evaluation inputs unless they are converted into train-ready image and label records. The \cmd{doctor import} report includes \path{reference_trainer.direct_train_ready} and \path{reference_trainer.train_ready_after_migration} so automation can distinguish direct training inputs from native sources that need migration first. The train-specific \cmd{doctor train-dataset} route performs the same check and emits next commands without launching training. It also recognizes explicit \cmd{classification}, \cmd{obb}, \cmd{depth}, and \cmd{pose6d} source families. Classification folders/manifests and OBB labels are recognized intake contracts for external training lanes, not direct RT-DETR reference-trainer inputs. Depth and pose6d are direct only when normalized bbox records include the required sidecars. Through the public CLI, \cmd{python3 -m yolozu train <config> --dataset-root <wrapper>} forwards the dataset root to the RT-DETR reference trainer.

For record-based training, \cmd{python3 -m yolozu train <config> --records-json <train.json> --val-records-json <val.json>} forwards explicit train and validation records to the reference trainer. Both files may be a JSON list of records or an object with an \path{images} list. Records may use either \path{image} or \path{image_path}; label boxes may use flat \path{cx}/\path{cy}/\path{w}/\path{h} fields or nested \path{bbox} fields. OBB records may include \path{angle} or \path{obb}; depth and pose6d records may carry sidecars at the record level or on individual labels. Use this path when the dataset cannot be cleanly represented as one YOLO-style split directory.

\begin{lstlisting}[language=bash]
python3 -m yolozu doctor train-dataset \
  --from auto \
  --dataset /path/to/native_dataset_root \
  --split val2017 \
  --output -
python3 -m yolozu doctor import \
  --dataset-from auto \
  --dataset /path/to/native_dataset_root \
  --split val2017 \
  --output -
python3 -m yolozu migrate dataset \
  --from auto \
  --dataset /path/to/native_dataset_root \
  --split val2017 \
  --output reports/native_dataset_wrapper \
  --force
python3 rtdetr_pose/tools/train_minimal.py \
  --dataset-root reports/native_dataset_wrapper \
  --config rtdetr_pose/configs/base.json \
  --max-steps 50
\end{lstlisting}

COCO keypoints roots can be selected explicitly when needed:
\begin{lstlisting}[language=bash]
python3 -m yolozu doctor train-dataset \
  --from coco-keypoints \
  --dataset /path/to/coco_keypoints_root \
  --split val2017 \
  --output -
python3 -m yolozu import dataset \
  --from coco-keypoints \
  --dataset /path/to/coco_keypoints_root \
  --split val2017 \
  --output reports/coco_keypoints_wrapper \
  --force
\end{lstlisting}

\section{The Run Contract: Rationale and implementation}
A standardized run directory contains:
\begin{itemize}
  \item Checkpoints (e.g., \path{runs/<run-id>/checkpoints/best.pt})
  \item Exported models (e.g., ONNX formats)
  \item Predictions and evaluation reports
  \item Calibration statistics artifacts (when explicitly enabled)
\end{itemize}

This layout makes these operations easier:
\begin{itemize}
  \item Reproducing experiment results
  \item Running parity checks
  \item Comparing runs across Git branches or compute nodes
\end{itemize}

Core rules of the run interface contract, documented here and enforced by the trainer:
\begin{itemize}
  \item Activation via the \cmd{--run-contract} or \cmd{--run-id <id>} flags.
  \item Run-interface-contract mode requires explicit configuration keys and rejects unrecognized keys to prevent silent configuration drift.
  \item A deterministic artifact layout under \path{runs/<run_id>/}:
    \begin{itemize}
      \item \path{checkpoints/last.pt}, \path{checkpoints/best.pt}
      \item \path{reports/train_metrics.jsonl}, \path{reports/val_metrics.jsonl}
      \item \path{reports/config_resolved.yaml}, \path{reports/run_meta.json}, \path{reports/training_summary.json}
    \end{itemize}
  \item Optional export/parity artifacts (when export/parity is enabled):
    \path{exports/model.onnx}, \path{exports/model.onnx.meta.json}, \path{reports/onnx_parity.json}.
  \item Full-state resumption is standardized using the \cmd{--resume} flag, defaulting to \path{runs/<run_id>/checkpoints/last.pt}.
\end{itemize}

\begin{figure}[h]
\centering
\resizebox{\linewidth}{!}{%
\begin{tikzpicture}[
  node distance=10mm,
  font=\small,
  box/.style={draw, rounded corners, align=left, inner sep=6pt},
  arrow/.style={-{Latex[length=2mm]}, thick}
]

\node[box] (root) {\textbf{runs/<run\_id>/}\\stable artifact root};

\node[box, below=14mm of root, xshift=-55mm] (ckpt) {\textbf{checkpoints/}\\\path{last.pt}\\\path{best.pt}};
\node[box, below=14mm of root] (reports) {\textbf{reports/}\\\path{train_metrics.jsonl}\\\path{val_metrics.jsonl}\\\path{config_resolved.yaml}\\\path{run_meta.json}};
\node[box, below=14mm of root, xshift=55mm] (exports) {\textbf{exports/} (optional)\\\path{model.onnx}\\\path{model.onnx.meta.json}};

\node[box, below=12mm of reports] (parity) {\textbf{parity/eval artifacts} (optional)\\\path{reports/onnx_parity.json}\\\path{reports/*.json}};

\draw[arrow] (root) -- (ckpt);
\draw[arrow] (root) -- (reports);
\draw[arrow] (root) -- (exports);
\draw[arrow] (reports) -- (parity);

\end{tikzpicture}
}
\caption{Run directory layout: fixed locations make resume, parity, and cross-branch comparisons reproducible.}
\end{figure}

\subsection{Artifact contract: required vs optional}
\begin{longtable}{@{}p{0.24\textwidth}p{0.26\textwidth}p{0.42\textwidth}@{}}
\toprule
\textbf{Class} & \textbf{Path} & \textbf{Condition} \\
\midrule
Required & \path{runs/<run_id>/checkpoints/last.pt} & Always in contracted runs \\
Required & \path{runs/<run_id>/reports/train_metrics.jsonl} & Always in contracted runs \\
Required & \path{runs/<run_id>/reports/val_metrics.jsonl} & Always in contracted runs \\
Required & \path{runs/<run_id>/reports/config_resolved.yaml} & Always in contracted runs \\
Required & \path{runs/<run_id>/reports/run_meta.json} & Always in contracted runs \\
Optional & \path{runs/<run_id>/exports/model.onnx} & Present when ONNX export is enabled \\
Optional & \path{runs/<run_id>/exports/model.onnx.meta.json} & Present when ONNX export metadata is emitted \\
Optional & \path{runs/<run_id>/reports/onnx_parity.json} & Present when ONNX parity check is executed \\
Optional & \path{runs/<run_id>/reports/training_summary.json} & Present as the backend-neutral top-level summary interface contract \\
\bottomrule
\end{longtable}


\noindent\textbf{Note (low-level trainer):} When invoking \cmd{python3 -m rtdetr_pose.train_minimal} directly, specifying \cmd{--run-id} implicitly activates run-contract mode and necessitates a \cmd{--config} argument (a YAML/JSON file such as \path{configs/examples/train_contract.yaml}). For rapid smoke tests bypassing the run contract, use explicit values such as \cmd{--run-dir runs/exp001} and \cmd{--onnx-out runs/exp001/model.onnx}.

\section{Configuration resolution and strictness}
The training entry point reads YAML/JSON configurations via \cmd{--config}, then applies explicit CLI overrides. The resolution hierarchy is:
\begin{center}
\textbf{CLI flags \(>\) config file \(>\) built-in defaults}
\end{center}

Configuration keys must correspond exactly to their argparse destination names (e.g., \cmd{--weight-decay} maps to \texttt{weight\_decay}). In run-contract mode (triggered by \texttt{run\_contract} or \texttt{run\_id}), any unrecognized keys are immediately rejected, thereby eliminating the risk of silent configuration drift.

\begin{figure}[h]
\centering
\resizebox{0.95\linewidth}{!}{%
\begin{tikzpicture}[
  node distance=10mm,
  font=\small,
  box/.style={draw, rounded corners, align=left, inner sep=6pt},
  arrow/.style={-{Latex[length=2mm]}, thick}
]

\node[box] (cli) {\textbf{CLI flags}\\explicit overrides\\(\cmd{--lr}, \cmd{--amp}, ...)};
\node[box, below=10mm of cli] (cfg) {\textbf{Config file}\\YAML/JSON via \cmd{--config}};
\node[box, below=10mm of cfg] (def) {\textbf{Built-in defaults}\\conservative baseline};

\node[box, right=28mm of cfg] (resolved) {\textbf{Resolved config}\\\path{config_resolved.yaml}\\(run artifact)};
\node[box, right=28mm of resolved] (strict) {\textbf{Strict mode}\\unknown keys fail\\(no silent drift)};

\draw[arrow] (cli) -- (resolved);
\draw[arrow] (cfg) -- (resolved);
\draw[arrow] (def) -- (resolved);
\draw[arrow] (resolved) -- (strict);

\end{tikzpicture}
}
\caption{Configuration resolution: the run contract enforces explicit keys and produces a resolved config artifact.}
\end{figure}

\section{YAML workflows: train / resume / test}
A unified project structure supports YAML-driven execution for both training and testing scenarios.

\begin{longtable}{@{}p{0.18\textwidth}p{0.30\textwidth}p{0.46\textwidth}@{}}
\toprule
\textbf{Workflow} & \textbf{Primary config} & \textbf{Reference command} \\
\midrule
Train (starter config) & \path{configs/examples/train_setting.yaml} & \cmd{yolozu train <train_setting.yaml>} \\
Train (run contract) & \path{configs/examples/train_contract.yaml} & \cmd{yolozu train <train_contract.yaml> --run-id <id>} \\
Resume (full state) & Same as contract & \cmd{yolozu train <train_contract.yaml> --run-id <id> --resume} \\
Test (scenario runner) & \path{configs/examples/test_setting.yaml} & \cmd{yolozu test <test_setting.yaml> [test_args...]} \\
\bottomrule
\end{longtable}

Example execution commands:
\begin{lstlisting}[language=bash]
yolozu train configs/examples/train_setting.yaml
yolozu train configs/examples/train_contract.yaml --run-id exp01
yolozu train configs/examples/train_contract.yaml --run-id exp01 --resume
yolozu test configs/examples/test_setting.yaml --adapter precomputed --predictions reports/predictions.json
\end{lstlisting}


\noindent For \cmd{yolozu test}, YAML keys are dynamically translated into \cmd{--key value} arguments (with booleans converted to flags), allowing subsequent CLI arguments to override them as necessary.

\section{Backend interface and orchestration}
The training platform now separates five concerns:
\begin{enumerate}
  \item a canonical \cmd{TrainConfig} projection,
  \item a backend registry,
  \item a shared \path{training_summary.json} style interface contract,
  \item a capability matrix, and
  \item a lightweight orchestration entrypoint.
\end{enumerate}

This means the reference trainer and the external lanes do not need to share
one internal loop, but they do share one top-level description of what ran and
what artifacts were promised.

For external lanes, YOLOZU now also standardizes one wrapper-level bundle under
\path{work_dir/}:
\begin{itemize}
  \item \path{dataset/}
  \item \path{configs/train_config_projection.json}
  \item \path{reports/training_summary.json}
  \item \path{reports/external_run_meta.json}
  \item \path{reports/launcher_plan.json}
  \item \path{reports/execution.json}
  \item \path{reports/resume_handoff.json}
  \item \path{reports/export_handoff.json}
  \item \path{reports/eval_handoff.json}
  \item \path{reports/parity_handoff.json}
  \item \path{reports/training_registry_entry.json}
\end{itemize}

This is not the same as forcing one backend-native checkpoint layout. It means
the wrapper-owned interface contract is now fixed even when the actual trainer
remains external.

The external training summary also records \path{next_steps}: copy-paste
follow-up commands for resume, export, evaluation, and parity. The handoff JSONs make
those next steps machine-readable, so orchestration and CI do not need to infer
backend-specific conventions from prose. In practice, this makes the training
report itself the hand-off note for the next operator step while still keeping
the resume/export/eval/parity interface contract explicit.

OpenCV DNN and ONNX Runtime do not appear in this training registry because
YOLOZU treats them as inference/export runtimes after training, not as
training backends in their own right.

The repo-side orchestration command is:
\begin{lstlisting}[language=bash]
yolozu train-orchestrate \
  --spec reports/train_orchestration_spec.json \
  --output reports/training_orchestration_report.json \
  --registry-out reports/training_registry.jsonl
\end{lstlisting}

When \cmd{--registry-out} is set, each executed training run appends one
\path{yolozu_training_registry_entry_v1} record to a shared JSONL file.
This is intentionally lightweight: it is an append-only experiment registry for
cross-backend auditing, not a full experiment tracking server.

Read next:
\begin{itemize}
  \item \path{docs/training_backend_interface.md}
  \item \path{docs/training_capability_matrix.md}
  \item \path{docs/training_orchestration.md}
\end{itemize}

\section{External framework finetune smoke matrix}
For user convenience, YOLOZU includes template configs and a unified smoke runner for external finetune entrypoints:
\begin{itemize}
  \item YOLO-family runtime
  \item Detectron2
  \item MMDetection
  \item MMPose
  \item MMSeg
  \item NVIDIA TAO
  \item RT-DETR (\path{rtdetr_pose})
\end{itemize}

Template paths:
\begin{itemize}
  \item \path{configs/examples/finetune_external/yolo_runtime_yolov8n_finetune_smoke.yaml}
  \item \path{configs/examples/finetune_external/mmdetection_finetune_smoke.py}
  \item \path{configs/examples/finetune_external/mmpose_finetune_smoke.py}
  \item \path{configs/examples/finetune_external/mmseg_finetune_smoke.py}
  \item \path{configs/examples/finetune_external/tao_finetune_smoke.yaml}
  \item \path{configs/examples/finetune_external/detectron2_finetune_smoke.yaml}
  \item \path{configs/examples/finetune_external/rtdetr_pose_finetune_smoke.yaml}
\end{itemize}

Unified smoke command (interface contract report):
\begin{lstlisting}[language=bash]
python3 tools/run_external_finetune_smoke.py \
  --dataset-root data/smoke \
  --split train \
  --output reports/external_finetune_smoke.json
\end{lstlisting}

Top-level Detectron2 training lane examples:
\begin{lstlisting}[language=bash]
python3 -m yolozu train \
  --external-backend detectron2 \
  configs/examples/finetune_external/detectron2_finetune_smoke.yaml \
  --dataset data/smoke \
  --split val \
  --task-family bbox \
  --dry-run \
  --output reports/train_external_detectron2_bbox.json

python3 -m yolozu train \
  --external-backend detectron2 \
  /path/to/mask_rcnn_config.yaml \
  --dataset /path/to/dataset \
  --split train \
  --task-family segmentation \
  --train-opt DATASETS.TRAIN "(\"my_seg_train\",)" \
  --dry-run \
  --output reports/train_external_detectron2_seg.json

python3 -m yolozu train \
  --external-backend detectron2 \
  /path/to/keypoint_rcnn_config.yaml \
  --dataset /path/to/dataset \
  --split train \
  --task-family keypoints \
  --train-opt DATASETS.TRAIN "(\"my_kpts_train\",)" \
  --dry-run \
  --output reports/train_external_detectron2_keypoints.json
\end{lstlisting}

Execute real training for selected frameworks:
\begin{lstlisting}[language=bash]
python3 tools/run_external_finetune_smoke.py \
  --dataset-root data/smoke \
  --split train \
  --non-dry-framework yolov \
  --non-dry-framework rtdetr \
  --epochs 1 --max-steps 1 --batch-size 2 --image-size 96 \
  --device cpu \
  --require-training-execution \
  --output reports/external_finetune_smoke.exec.json
\end{lstlisting}

MMDetection/Detectron2 can be wired to external launchers via
\cmd{--mmdet-train-script} and \cmd{--detectron2-train-script}.
RT-DETR non-dry dependency failures are explicit in the report
(\texttt{failure\_code=E\_DEP\_TORCH\_MISSING}).
When MMDetection/Detectron2 projection dependencies are absent, train-path audit can still continue with external launchers; the report records \texttt{projection\_error} and \texttt{train\_path\_audited}.

GPU check (machine.dev style) example:
\begin{lstlisting}[language=bash]
python3 tools/run_external_finetune_smoke.py \
  --dataset-root data/smoke \
  --split train \
  --non-dry-framework rtdetr \
  --device cuda \
  --epochs 1 --max-steps 1 --batch-size 2 --image-size 96 \
  --require-training-execution \
  --output reports/external_finetune_smoke.machine_dev.json
\end{lstlisting}

\section{Solver (optimizer) parameters and defaults}
The supported optimizers are \texttt{adamw} (the default) and \texttt{sgd}.

\begin{longtable}{@{}p{0.40\textwidth}p{0.10\textwidth}p{0.12\textwidth}p{0.30\textwidth}@{}}
\toprule
\textbf{CLI / config key} & \textbf{Type} & \textbf{Default} & \textbf{Notes} \\
\midrule
\cmd{--optimizer} / \texttt{optimizer} & str & \texttt{adamw} & \texttt{adamw}, \texttt{sgd} \\
\cmd{--lr} / \texttt{lr} & float & \texttt{1e-4} & Base learning rate \\
\cmd{--weight-decay} / \texttt{weight\_decay} & float & \texttt{0.01} & Base weight decay \\
\cmd{--momentum} / \texttt{momentum} & float & \texttt{0.9} & Applicable to SGD only \\
\cmd{--nesterov} / \texttt{nesterov} & bool & \texttt{false} & Applicable to SGD only \\
\cmd{--use-param-groups} / \texttt{use\_param\_groups} & bool & \texttt{false} & Enables separate backbone/head parameter groups \\
\cmd{--backbone-lr-mult} / \texttt{backbone\_lr\_mult} & float & \texttt{1.0} & Learning rate multiplier for the backbone \\
\cmd{--head-lr-mult} / \texttt{head\_lr\_mult} & float & \texttt{1.0} & Learning rate multiplier for the head \\
\cmd{--backbone-wd-mult} / \texttt{backbone\_wd\_mult} & float & \texttt{1.0} & Weight decay multiplier for the backbone \\
\cmd{--head-wd-mult} / \texttt{head\_wd\_mult} & float & \texttt{1.0} & Weight decay multiplier for the head \\
\cmd{--wd-exclude-bias} / \texttt{wd\_exclude\_bias} & bool & \texttt{true} & Sets bias decay to 0 \\
\cmd{--wd-exclude-norm} / \texttt{wd\_exclude\_norm} & bool & \texttt{true} & Sets normalization decay to 0 \\
\bottomrule
\end{longtable}

\section{LR scheduling parameters and defaults}
Supported learning rate scheduler configurations include \texttt{none} (default), \texttt{cosine}, \texttt{onecycle}, and \texttt{multistep}.

\begin{longtable}{@{}p{0.40\textwidth}p{0.10\textwidth}p{0.12\textwidth}p{0.30\textwidth}@{}}
\toprule
\textbf{CLI / config key} & \textbf{Type} & \textbf{Default} & \textbf{Notes} \\
\midrule
\cmd{--scheduler} / \texttt{scheduler} & str & \texttt{none} & \texttt{none}, \texttt{cosine}, \texttt{onecycle}, \texttt{multistep} \\
\cmd{--min-lr} / \texttt{min\_lr} & float & \texttt{0.0} & Corresponds to cosine \texttt{eta\_min} \\
\cmd{--scheduler-milestones} / \texttt{scheduler\_milestones} & str/list & \texttt{""} & Comma-separated list for multistep scheduling \\
\cmd{--scheduler-gamma} / \texttt{scheduler\_gamma} & float & \texttt{0.1} & Decay factor for multistep scheduling \\
\cmd{--lr-warmup-steps} / \texttt{lr\_warmup\_steps} & int & \texttt{0} & Number of linear warmup steps \\
\cmd{--lr-warmup-init} / \texttt{lr\_warmup\_init} & float & \texttt{0.0} & Initial learning rate during warmup \\
\bottomrule
\end{longtable}


\noindent\textbf{Note:} The \texttt{linear} scheduler is not supported in the current codebase.

\section{LoRA / QLoRA parameters and defaults}
Low-Rank Adaptation (LoRA) is activated when \texttt{lora\_r > 0}.
For the plain-language intuition and the difference between LoRA and QLoRA, see Chapter~14 before tuning these knobs.

\paragraph{What these words mean in plain language.}
\begin{itemize}
  \item \textbf{LoRA} stands for \textbf{Low-Rank Adaptation}. Instead of updating every weight in a large model, LoRA adds a small trainable adapter on top of selected layers. In practice, this means ``keep most of the base model fixed, and learn a compact correction.''
  \item \textbf{Why people use LoRA}: it reduces memory cost, usually trains faster than full-model finetuning, and makes it easier to keep one base checkpoint while testing several small adaptation variants.
  \item \textbf{\texttt{lora\_r}} is the adapter rank. A larger rank means a more expressive adapter, but also more trainable parameters.
  \item \textbf{\texttt{lora\_alpha}} is a scaling factor for the adapter contribution. A simple mental model is: rank controls adapter capacity, alpha controls how strongly the adapter is allowed to matter.
  \item \textbf{QLoRA} means ``LoRA on top of a quantized base model.'' In this repository, \texttt{--qlora} keeps the trainable LoRA adapters, but stores the frozen base model in a lower-precision quantized form to reduce memory further.
  \item \textbf{Why QLoRA is stricter}: because quantization and adapter training interact. That is why \yolozu{} enforces the safe combination explicitly instead of pretending every LoRA run is also a safe QLoRA run.
\end{itemize}

\paragraph{A simple comparison.}
\begin{itemize}
  \item \textbf{Full finetuning}: move almost every parameter.
  \item \textbf{LoRA finetuning}: keep the base model mostly fixed and train small adapters.
  \item \textbf{QLoRA finetuning}: same adapter idea as LoRA, but the frozen base model is additionally quantized to save memory.
\end{itemize}

\begin{longtable}{@{}p{0.40\textwidth}p{0.10\textwidth}p{0.12\textwidth}p{0.30\textwidth}@{}}
\toprule
\textbf{CLI / config key} & \textbf{Type} & \textbf{Default} & \textbf{Notes} \\
\midrule
\cmd{--lora-r} / \texttt{lora\_r} & int & \texttt{0} & Values \texttt{>0} enable LoRA \\
\cmd{--lora-alpha} / \texttt{lora\_alpha} & float/null & \texttt{null} & A null value implies \(\alpha=r\) \\
\cmd{--lora-dropout} / \texttt{lora\_dropout} & float & \texttt{0.0} & Specifies LoRA input dropout \\
\cmd{--lora-target} / \texttt{lora\_target} & str & \texttt{head} & Options: \texttt{head}, \texttt{all\_linear}, \texttt{all\_conv1x1}, \texttt{all\_linear\_conv1x1} \\
\cmd{--lora-freeze-base} / \texttt{lora\_freeze\_base} & bool & \texttt{true} & Restricts training exclusively to LoRA parameters \\
\cmd{--lora-train-bias} / \texttt{lora\_train\_bias} & str & \texttt{none} & Options: \texttt{none}, \texttt{all} \\
\cmd{--torchao-quant} / \texttt{torchao\_quant} & str & \texttt{none} & Options: \texttt{none}, \texttt{int8wo}, \texttt{int4wo} \\
\cmd{--torchao-required} / \texttt{torchao\_required} & bool & \texttt{false} & Triggers a failure if the quantization backend is absent \\
\cmd{--qlora} / \texttt{qlora} & bool & \texttt{false} & Requires LoRA; strictly enforces int4wo quantization and base freezing \\
\bottomrule
\end{longtable}

\section{Default configuration list (practical minimum)}
The canonical default configurations are maintained within:
\begin{itemize}
  \item \path{configs/examples/train_setting.yaml}
  \item \path{configs/examples/train_contract.yaml}
  \item \path{configs/examples/test_setting.yaml}
\end{itemize}

Frequently modified default parameters include:
\begin{longtable}{@{}p{0.34\textwidth}p{0.20\textwidth}p{0.40\textwidth}@{}}
\toprule
\textbf{Config key} & \textbf{Default} & \textbf{Meaning} \\
\midrule
\texttt{optimizer} & \texttt{adamw} & Optimizer family \\
\texttt{lr} & \texttt{1e-4} & Base learning rate \\
\texttt{weight\_decay} & \texttt{0.01} & Base weight decay \\
\texttt{scheduler} & \texttt{none} & Learning rate scheduler mode \\
\texttt{lr\_warmup\_steps} & \texttt{0} & Warmup is disabled by default \\
\texttt{lora\_r} & \texttt{0} & LoRA is disabled by default \\
\texttt{gradient\_accumulation\_steps} & \texttt{1} & Gradient accumulation is disabled \\
\texttt{clip\_grad\_norm} & \texttt{0.0} & Gradient clipping is disabled \\
\texttt{use\_ema} & \texttt{false} & Exponential Moving Average (EMA) is disabled \\
\texttt{amp} & \texttt{none} & Automatic Mixed Precision (AMP) is disabled \\
\texttt{metrics\_jsonl} & \path{reports/train_metrics.jsonl} & Output stream for training metrics \\
\texttt{metrics\_csv} & \path{reports/train_metrics.csv} & Tabular format for metrics \\
\texttt{export\_onnx} & \texttt{true} & ONNX export is enabled \\
\texttt{onnx\_out} & \path{reports/rtdetr_pose.onnx} & Designated ONNX output path \\
\bottomrule
\end{longtable}

\section{Smoke training}
For a repository-wide offline smoke, use:
\begin{lstlisting}[language=bash]
bash scripts/smoke.sh
\end{lstlisting}

For training-entrypoint coverage across external frameworks, use the dedicated
interface-contract smoke runner:
\begin{lstlisting}[language=bash]
python3 tools/run_external_finetune_smoke.py \
  --dataset-root data/smoke \
  --split train \
  --output reports/external_finetune_smoke.json
\end{lstlisting}

The first command exercises the main offline smoke lane. The second focuses on
finetune-entrypoint audit for YOLOv/MMDetection/Detectron2/RT-DETR style
training surfaces and writes a machine-readable report.

\section{Depth mode in training (current implementation)}
The in-repo RT-DETR Pose trainer incorporates optional depth integration, configured with a conservative default:
\begin{itemize}
  \item \cmd{--depth-mode none} (default): Executes the baseline pathway with depth processing disabled.
  \item \cmd{--depth-mode sidecar}: Ingests per-image depth sidecars (\cmd{depth_path}/\cmd{depth}) and propagates the \cmd{depth_valid} flag.
  \item \cmd{--depth-mode fuse_mid}: Implements lightweight depth fusion subsequent to the projector (external to the backbone boundary).
\end{itemize}

Safety controls:
\begin{itemize}
  \item \cmd{--depth-unit \{unspecified,relative,metric\}} (default: \cmd{unspecified}).
  \item Absolute depth matcher costs are strictly permitted only in metric mode; non-metric modes automatically disable \cmd{cost_z}/\cmd{cost_t} for safety.
  \item \cmd{--depth-scale} applies a scalar multiplier to sidecar depth values.
  \item \cmd{--depth-dropout} regulates the modality dropout probability during \cmd{fuse_mid} training.
\end{itemize}

\section{Contracted run example}
\begin{lstlisting}[language=bash]
yolozu train configs/examples/train_contract.yaml --run-id exp01

# Resume training from runs/exp01/checkpoints/last.pt
yolozu train configs/examples/train_contract.yaml --run-id exp01 --resume
\end{lstlisting}

\section{Backbone modification (adapter-scoped architecture switch)}
Backbone-specific configurations are not primarily governed by solver flags; rather, they are derived from the model configuration file specified by \texttt{model\_config} in the training YAML.
This section applies to the in-repo \path{rtdetr_pose} adapter boundary.

Backbone notes and examples are maintained in \path{docs/backbones.md}.

The backbone output contract strictly mandates three feature maps: \texttt{[P3, P4, P5]} with corresponding strides of \texttt{[8, 16, 32]}. Each level is independently projected to the transformer's \texttt{d\_model} dimension via 1x1 convolutions, then processed by the RT-DETR-style FPN/PAN + hybrid encoder path.

Supported backbone architectures in the current codebase include:
\texttt{cspresnet}, \texttt{tiny\_cnn}, \texttt{cspdarknet\_s}, \texttt{resnet50}, and \texttt{convnext\_tiny}.

While legacy compatibility fields (e.g., \texttt{backbone\_name}, \texttt{backbone\_channels}) remain functional, the strongly preferred convention is to use \path{model.backbone.*} together with \path{model.projector.d_model}.

\begin{longtable}{@{}p{0.28\textwidth}p{0.24\textwidth}p{0.42\textwidth}@{}}
\toprule
\textbf{Location} & \textbf{Key} & \textbf{Role} \\
\midrule
Train YAML & \texttt{model\_config} & References the model JSON (e.g., \path{rtdetr_pose/configs/base.json}) \\
Model JSON (\texttt{model.backbone.*}) & \texttt{name} & Specifies the backbone family (e.g., \texttt{cspresnet}, \texttt{cspdarknet\_s}, \texttt{resnet50}, \texttt{convnext\_tiny}) \\
Model JSON (\texttt{model.backbone.*}) & \texttt{norm} & Defines the normalization type (\texttt{bn|syncbn|frozenbn|gn}) \\
Model JSON (\texttt{model.backbone.*}) & \texttt{args} & Optional, backbone-specific constructor arguments \\
Model JSON (\texttt{model.projector.*}) & \texttt{d\_model} & Sets the projection/output channel size for transformer input \\
Solver YAML/CLI & \texttt{backbone\_lr\_mult} & Learning rate multiplier applied to backbone parameters \\
Solver YAML/CLI & \texttt{backbone\_wd\_mult} & Weight decay multiplier applied to backbone parameters \\
\bottomrule
\end{longtable}

Normalization guidelines:
\begin{itemize}
  \item \texttt{bn}: The standard default for single-node training.
  \item \texttt{syncbn}: Highly recommended for multi-GPU environments utilizing small per-rank batch sizes.
  \item \texttt{frozenbn}: Provides stability when batch statistics must remain static.
  \item \texttt{gn}: A robust, batch-size-independent alternative.
\end{itemize}

Recommended workflow for backbone modification:
\begin{enumerate}
  \item Duplicate \path{rtdetr_pose/configs/base.json} to create a project-local model JSON file.
  \item Modify \texttt{model.backbone.name} and \texttt{model.backbone.args} as required.
  \item Configure \texttt{model.projector.d\_model} to match the intended transformer hidden size.
  \item Update the training YAML's \texttt{model\_config} to reference the newly created JSON file.
  \item Initiate a smoke test with minimal settings (e.g., reduced \texttt{image\_size} and \texttt{max\_steps}) before scaling up.
\end{enumerate}

Example implementation:
\begin{lstlisting}[language=bash]
cp rtdetr_pose/configs/base.json rtdetr_pose/configs/my_backbone.json
# Edit model.backbone.name, model.backbone.args, and model.projector.d_model accordingly

yolozu train configs/examples/train_setting.yaml   --model-config rtdetr_pose/configs/my_backbone.json   --use-param-groups   --backbone-lr-mult 0.1
\end{lstlisting}

\section{PyTorch utility wrappers}
The \texttt{yolozu.training.torch\_utils} module provides thin wrappers around core PyTorch APIs for common training and inference tasks:

\begin{itemize}
  \item \texttt{amp\_inference\_context()} --- wraps \texttt{torch.amp.autocast} for mixed-precision inference.
  \item \texttt{compile\_model()} --- wraps \texttt{torch.compile} with configurable backend/mode.
  \item \texttt{profile\_callable()} --- wraps \texttt{torch.profiler} and returns structured event summaries.
  \item \texttt{model\_info()} --- reports parameter counts, dtype breakdown, device, and buffer stats.
  \item \texttt{auto\_device()} --- auto-detects the best available device (CUDA $\to$ MPS $\to$ CPU).
  \item \texttt{seed\_everything()} --- sets all random seeds (Python, NumPy, PyTorch) for reproducibility.
\end{itemize}

\begin{lstlisting}[language=Python]
from yolozu.training.torch_utils import (
    amp_inference_context, compile_model, model_info,
    auto_device, seed_everything,
)

seed_everything(42)
device = auto_device()
model = model.to(device)
compiled = compile_model(model, mode="reduce-overhead")
print(model_info(compiled))

with amp_inference_context(device.type):
    output = compiled(images)
\end{lstlisting}

\section{AMP utilities for training loops}
The \texttt{yolozu.training.amp\_utils} module provides a thin wrapper around \texttt{torch.amp} that standardizes Automatic Mixed Precision setup across SDFT distillation, TTA/TTT loops, and custom training scripts.

\begin{itemize}
  \item \texttt{amp\_available()} --- returns \texttt{True} when \texttt{torch.amp.autocast} is present (PyTorch 2.x).
  \item \texttt{make\_amp\_context(device\_type, dtype, enabled)} --- creates a paired \texttt{(autocast\_factory, GradScaler|None)} tuple. A \texttt{GradScaler} is only returned for \texttt{float16} on CUDA; for \texttt{bfloat16}, MPS, or CPU the scaler is \texttt{None}.
\end{itemize}

\begin{lstlisting}[language=Python]
from yolozu.training.amp_utils import make_amp_context

autocast_ctx, scaler = make_amp_context(device_type="cuda", dtype="float16")
with autocast_ctx():
    loss = model(inputs)
if scaler is not None:
    scaler.scale(loss).backward()
    scaler.step(optimizer)
    scaler.update()
else:
    loss.backward()
    optimizer.step()
\end{lstlisting}

\section{Torchvision v2 transforms bridge}
The \texttt{yolozu.training.transforms\_bridge} module provides composable augmentation pipelines using \texttt{torchvision.transforms.v2} (stable since torchvision 0.17). These pipelines jointly transform images, bounding boxes, and keypoints in a single call.

\begin{itemize}
  \item \texttt{transforms\_v2\_available()} --- checks whether \texttt{torchvision.transforms.v2} is importable.
  \item \texttt{build\_detection\_transforms(size, hflip\_prob, color\_jitter, scale\_range, ...)} --- constructs a training pipeline (random resize, horizontal flip, colour jitter, normalize).
  \item \texttt{build\_eval\_transforms(size, ...)} --- constructs an evaluation pipeline (resize + normalize only).
\end{itemize}

Both builders insert \texttt{v2.ToImage()} before \texttt{v2.Normalize()} so that raw PIL inputs are converted to tensors first, avoiding the common ``Normalize requires tensor'' error.

\begin{lstlisting}[language=Python]
from yolozu.training.transforms_bridge import (
    build_detection_transforms, build_eval_transforms,
)

train_tfm = build_detection_transforms(size=(640, 640), hflip_prob=0.5)
eval_tfm  = build_eval_transforms(size=(640, 640))
\end{lstlisting}
